\documentclass{article}
\usepackage{geometry}
\geometry{a4paper,scale=0.8}
\usepackage[utf8]{inputenc}

\usepackage{amsmath, amssymb}
\usepackage{amsthm}
\usepackage{enumitem}

\newtheoremstyle{breakstyle}
    {5pt}
    {10pt}
    % {\normalfont}
    {\itshape}
    {0pt}
    {\bfseries}
    {\newline\indent}
    {0pt}
    {\thmname{#1}\thmnumber{ #2}\thmnote{ \textbf{[#3]}}}
\theoremstyle{breakstyle}
\newtheorem{exercise}{Exercise}[section]
\newtheorem{remark}{Remark}[section]

\newenvironment{solution}
  {\begin{proof}[Solution.]\renewcommand{\qedsymbol}{}}
  {\end{proof}}

% \let\ker\relax
% \DeclareMathOperator{\ker}{Ker}
\let\det\relax
\DeclareMathOperator{\det}{Det}

\title{Chapter 4 - Affine reflection groups}
\author{Flandre}
\date{\today}

\begin{document}
\maketitle

\setcounter{section}{2}
\section{Alcoves}

\begin{exercise}[Trivial]
An alcove $A$ consists of all $\lambda \in V$ satisfying strict inequalities
$k_{\alpha} < (\lambda, \alpha) < k_{\alpha} + 1$ , where $\alpha$ runs over
$\Phi^{+}$ and $k_{\alpha} \in \mathbf{Z}$ . Its \textbf{upper closure} consists
of those $\lambda$ satisfying the inequalities obtained by replacing the second
$<$ by $\leq$ in each case. Prove that each $\lambda \in V$ lies in the upper
closure of a unique alcove.

\begin{proof}
    \emph{(Check the definition of those walls.)}
\end{proof}

\end{exercise}

\end{document}
